\documentclass{internal-report}
\usepackage[UTF8]{ctex}
\usepackage{amsmath, amssymb, amsthm}
\usepackage{booktabs}
\usepackage{graphicx}
\usepackage{hyperref}
\usepackage{float}

\title{问题 3 内部报告：多因素下的 BMI 分组与最佳 NIPT 时点}
\author{MathModel AI}
\date{2026-07-19}
% STATUS: done
% LAST_MODIFIED: 2026-07-19

\begin{document}

\maketitle

\section{问题重述与分析}

\subsection{题目要求与增量定位}

问题 3 是问题 2 的直接延续：在已知 BMI 是达标时间主要因素的前提下，进一步追问——是否还有其他因素（身高、体重、年龄、怀孕次数等）具有独立的预测贡献，以至于需要据此调整 BMI 分组方案？同时，需以达标比例（而非仅达标时间中位数）验证各组推荐时点的有效性，并考察检测误差的扰动。

因此本问题的增量贡献不在于提出新分组，而在于\textbf{穷举检验所有候选因素、排除虚假信号、为问题 2 的分组方案提供独立验证}。

\subsection{待检验的候选因素}

根据赛题提示和数据可用性，候选因素包括：身高、体重、年龄、怀孕次数、生产次数。辅助生育个体（IVF 8 例 + IUI 9 例，共 17 人）中仅 4 人有首次达标记录，样本量不足且临床管理路径特殊（全程生殖中心增强监护），不纳入多因素分析。

\subsection{检验策略}

对 260 位有首次达标记录的孕妇（排除辅助生育后 256 人），以首次达标孕周为因变量，逐个检验各候选因素的边际解释力（$\Delta R^2$）和统计显著性（$p$ 值）。若某因素的效应量足以在临床层面上改变推荐时点（$\geq 2$ 周），则认为需要对分组方案做相应调整；否则维持问题 2 的三层分组。

\section{多因素回归：哪些变量有实际贡献？}

\subsection{方法}

对 260 位有首次达标记录的孕妇（排除辅助生育个体后为 256 人；原始数据含 IVF 8 例 + IUI 9 例共 17 人，其中仅 4 人有首次达标记录），以首次达标孕周为因变量，将 BMI、身高、体重、年龄、怀孕次数、生产次数等变量纳入回归。

\subsection{结果}

\begin{table}[H]
\centering
\caption{多因素回归结果（排除辅助生育个体，n=256）}
\begin{tabular}{lrrr}
\toprule
变量 & 单变量 $R^2$ & 精简全模型 $p$ 值 & 评注 \\
\midrule
BMI     & 0.027 & 0.010 & 主要预测因子 \\
体重     & 0.045 & —（与 BMI 共线，未纳入精简模型） & \\
身高     & 0.017 & —（与 BMI 共线，未纳入精简模型） & \\
年龄     & 0.001 & 0.710 & 无显著贡献 \\
怀孕次数   & 0.001 & 0.788 & 无显著贡献 \\
生产次数   & 0.000 & 0.846 & 无显著贡献 \\
\bottomrule
\multicolumn{4}{l}{\footnotesize 精简全模型（BMI + 年龄 + 怀孕/生产次数）$R^2=0.028$；含体重/身高的全模型 $R^2=0.095$} \\
\multicolumn{4}{l}{\footnotesize 辅助生育个体（IVF 8 例 + IUI 9 例，共 17 人，其中 4 人有首次达标记录）因样本量不足且临床路径特殊已排除。} \\
\end{tabular}
\end{table}

\textbf{关键发现}：

\begin{enumerate}
    \item \textbf{BMI 是唯一显著预测因子。} 排除辅助生育个体后，年龄、怀孕次数、生产次数的 $p$ 值均 $>0.6$，增量 $R^2 < 0.001$——无证据支持任何额外变量纳入分组决策。
    \item \textbf{身高的贡献是共线带来的。} 身高与 BMI 隐式相关（BMI = 体重 / 身高$^2$），移除体重/BMI 后身高的独立贡献可忽略。
    \item \textbf{多因素模型不改变分组方案。} 年龄、怀孕次数、生产次数的增量 $R^2$ 均 $<0.001$（$p>0.6$），辅助生育因样本量不足且临床管理路径特殊已排除。数据集内所有候选变量均已穷举检验，无证据支持任何变量改变 BMI 分组策略。
\end{enumerate}

\subsection{结论：多因素无实质贡献，分组方案与问题 2 一致}

多因素回归证实了问题 2 的结论：BMI（及共线体重）是唯一有实际贡献的因素。排除辅助生育个体后，年龄、怀孕次数、生产次数均无显著效应（$p > 0.6$）。因此，问题 3 的最优 BMI 分组及最佳 NIPT 时点\textbf{与问题 2 保持一致}。

\subsection{GC 含量同样无贡献}

问题 1 的模型 A 和问题 2 的达标率对比已确认：GC 含量对 Y 染色体浓度无显著影响（$t=1.34, p=0.18$），GC 正常组与偏低组的达标率曲线无差异（组间差 $<2$ 个百分点）。其机制在于 Y 浓度 = Y reads / 总 reads 的比值形式天然约消了 GC 偏差的同向影响。因此，GC 含量不在问题 3 的多因素候选之列——其生物学特性使其不具备作为预测因子的条件，问题 1/2 的结论在此同样有效，不再重复检验。

\section{达标比例验证}

使用与问题 2 相同的早测低代价策略（$p_0=0.5$），对各组在建议时点的 Y $\geq$ 4\% 比例进行直接验证。

\begin{table}[H]
\centering
\caption{各组达标比例验证}
\begin{tabular}{lcccc}
\toprule
BMI 组 & 建议时点 & 该时点达标比例 & 该时点+补检达标比例 & 最终失败率 \\
\midrule
$[20,28)$   & 12 周 & 80\% & 95\% & $\leq 5.0\%$ \\
$[28,32)$   & 12 周 & 83\% & 100\% & 0.3\% \\
$[32,36)$  & 12 周 & 84\% & 99\% & 0.6\% \\
$[36,40)$ & 14 周 & 75\% & 92\% & 8.3\% \\
40+    & 14 周预检 & 33\% & 80\%→92\% & 8.0\% \\
\bottomrule
\end{tabular}
\end{table}

\begin{itemize}
    \item \textbf{轻体重组（BMI $<$ 36）}：单次检测达标比例 $\geq 83\%$，达标比例充分。补检一次后几乎全覆盖。
    \item \textbf{中体重组（BMI $[36,40)$）}：14 周首次达标比例 75\%，加 17--18 周复检后升至 92\%。达标比例为可接受水平。
    \item \textbf{重体组（40+）}：单次达标比例仅 33\%，需三次检测才升到 92\%。达标比例是该组主要的临床关注点。
\end{itemize}

\textbf{结论}：各组的达标比例在建议时点处于可接受范围。BMI 不是影响达标比例的唯一变量——即使各组已经是最优时点，仍有残差比例不达标，这源于个体生物学差异，与分组策略无关。

\section{检测误差对达标比例和时点的影响}

\subsection{四项误差来源}

| 误差来源 | 量级估计 | 对建议时点的影响 |
|---------|---------|----------------|
| Y 浓度测量误差（测序） | $\pm 0.5\%$ | 轻体重组无影响；40+ 组可能延迟 1--2 周 |
| 孕周测量误差（末次月经） | $\pm 1$ 周 | 所有组时点浮动 $\pm 1$ 周（临床预设中已包含） |
| 个体间基线变异（生物学） | CV 17--39\% | **最主要**：BMI $<$ 36 组 15\% 的个体达标延后 |
| $p_0$ 选择误差 | $\pm 0.1$ | 轻体重组无影响；40+ 组显著敏感 |

\subsection{系统性误差}

当 Y 浓度阈值因测序误差而被有效提高至 4.5\% 时：

\begin{itemize}
    \item $[28,36)$：达标比例从 83--84\% 轻微降至 78--80\%。缓冲余量充足。
    \item $[36,40)$：达标比例从 75\% 降至约 65\%——时点建议不变，但达标把握期望降低。
    \item 40+：20 周达标比例从 67\% 降至 50\%——系统性误差在此组放大。两组数据叠加后，（考虑多因素）无法为该组提供 80\% 以上把握的单一时点。
\end{itemize}

\subsection{误差影响总结}

检测误差对分组方案和最佳时点均不产生实质性修正——轻体组缓冲充分，中体组 14 周仍为最优，重度组数据本身是比测量误差更大的不确定源。测量误差通过 $p_0$ 的灵敏度分析和 $\pm 0.5\%$ 的系统偏移两个维度已充分覆盖。

\section{最终方案}

问题 3 的最优分组及最佳时点\textbf{与问题 2 保持一致}。三个层级：

$$
\begin{array}{cccc}
\hline
\text{层级} & \text{BMI 范围} & \text{建议检测方案} & \text{达标比例} \\
\hline
\text{标准层} & [20, 36) & \mathbf{12\text{ 周单检}} \rightarrow \text{若不达标 14--15 周} & 83--94\% \\
\text{中期层} & [36, 40) & \mathbf{14\text{ 周首检}} \rightarrow \text{若不达标 17--18 周} & 75\% \\
\text{预检层} & 40+ & \mathbf{14\text{ 周预检}} \rightarrow \text{18--20 周复检 + 确认} & 33\% \\
\hline
\end{array}
$$

\textbf{理由}：多因素回归证实 BMI（及共线体重）是唯一显著贡献的变量。排除辅助生育个体后，年龄、怀孕次数、生产次数的 $p$ 值均 $>0.6$，无证据支持纳入分组决策。辅助生育孕妇全程处于增强监护，NIPT 时点由主治医生个体化判断，不适用普查式分组方案。各组的达标比例已验证，检测误差不作修正。

\end{document}
